\bookchapter{제13장: 유틸리티 앱 추천}

\chapterintro{이 장은 생산성, 건강, 여행 도구를 엄선하고, 유용한 시스템·크로스 기기 기능을 덧붙입니다. 청탁(정리)·스팸 차단, 휴대폰 기본 차단 스위치, 날씨·지진 경보, Sunlogin/ToDesk 원격 제어, Trae 등 AI 오피스, 동적 배경 화면, LocalSend 등 LAN 고속 전송, 크로스 기기 클립보드, Aegis/Microsoft Authenticator OTP, TickTick 일정, 종이 문서 스캔 등.}

\booksection{생산성 도구: 시간과 집중력 관리}

생산성 도구는 디지털 생활의 “뼈대”입니다. 흩어진 아이디어, 할 일, 자료를 말끔하게 정리해 줍니다. 아래는 장면별 추천입니다.

\booksubsection{할 일 목록(작업 관리)}

\begin{itemize}
  \item 디다 리스트(TickTick) 이런 분께 좋아요: 업무, 생활, 공부 할 일과 일정을 한곳에서 통합 관리하고 싶은 분. 이 절에서 핵심 추천하는 앱 중 하나예요. 핵심 장점: 휴대폰/PC/태블릿/웹 어디서나 동기화되고, 리스트 + 캘린더 + 뽀모도로를 한 번에 쓸 수 있어요. '내일 오후 3시에 회의', '매주 월요일에 주간 보고 제출'처럼 자연어로 입력하면 시간과 반복 규칙을 알아서 인식하고, 하위 작업, 태그, 우선순위, 첨부 파일, 공동 작업 리스트도 지원해요. '일정 관리 + 할 일 체크 + 하루 돌아보기'를 하나로 관리하기에 딱이에요.
  \item TickTick을 더 편하게 쓰는 법: 수집함: 일단 전부 "오늘/수집함"에 넣고 저녁에 업무·가정·학습 목록으로 분류하세요.
  \item 캘린더 뷰: 회의와 마감일을 달력에 끌어다 놓으면 충돌이 한눈에 보입니다. 할 일 목록만 보는 것보다 직관적입니다.
  \item 반복 작업: 약 복용, 요금 납부, 운동은 규칙 작업으로 등록해 두면 매번 입력할 필요가 없습니다.
  \item 뽀모도로: 작업 위에서 바로 집중 모드를 켜고 이번 주에 "정말 끝낸 것"을 통계로 확인하세요.
  \item 휴대폰 기능 연동: 음성 비서에게 "\textasciitilde{}해줘"라고 하면 시스템 알림으로 들어갑니다. 중요한 마감일은 시스템 캘린더에도 등록해 이중 보완하세요.
\end{itemize}

Todoist 대상: 미니멀·크로스 플랫폼을 선호하는 사용자. 핵심 강점: 절제된 깔끔한 UI, 명확한 프로젝트 계층(작업하위 작업하위하위 작업), karma 포인트로 습관 형성 유도.

Microsoft To Do 대상: 오피스 사용자, Outlook 사용자. 핵심 강점: 완전 무료, 광고 없음. "내 일" 기능이 매일 비워 심리적 부담을 줄여 주고, 마이크로소프트 생태계와 깊게 연동됩니다.

시스템 메모/미리 알림 대상: 타사 앱을 깔기 싫은 가벼운 사용자. 핵심 강점: 네이티브 체험, 제로 지연. Siri/시리 음성으로 바로 추가 가능, 시스템 위젯과 연동.

\booksubsection{뽀모도로(집중력 훈련)}

\begin{itemize}
  \item Forest 대상: 손을 못 참고 자꾸 휴대폰을 보는 사람. 핵심 강점: 집중하면 나무를 심고, 미리 앱을 떠나면 나무가 말라 죽습니다. "가상 생명"으로 긍정적 피드백을 만들고 친구와 함께 심을 수도 있습니다.
  \item Tide 대상: 백색 소음으로 집중/수면이 필요한 사람. 핵심 강점: 자연 사운드(비, 파도, 숲) + 뽀모도로 + 명상 3-in-1. 미니멀하고 치유적인 UI.
  \item 원칙: 클래식 25분 집중 + 5분 휴식. 상태에 따라 50/10 또는 90/20으로 조정할 수 있습니다.
\end{itemize}

\booksubsection{캘린더(시간 시각화)}

\begin{itemize}
  \item 시스템 캘린더 대상: 모든 사용자. 핵심 강점: 네이티브 푸시가 가장 안정적이고 배터리 소모가 극히 낮습니다. Siri/시리 음성으로도 추가 가능.
  \item Google Calendar 대상: 다중 기기 사용자, 일정 공유가 필요한 사람. 핵심 강점: "목표" 기능이 운동/학습 시간을 자동 배치하고, Gmail에서 메일 속 일정을 스마트하게 인식합니다.
\end{itemize}

\booksubsection{노트(지식 축적)}

\begin{itemize}
  \item Notion 대상: 노트·할 일·데이터베이스·위키를 통합 관리하고 싶은 사람. 학생, 프리랜서, 팀 협업. 핵심 강점: 올인원 지식 베이스. Block 모듈형 편집으로 개인 홈페이지, 독서 목록, 프로젝트 보드를 만들 수 있고 템플릿 생태계가 풍부합니다.
  \item Obsidian 대상: 깊이 있는 지식 관리자, 연구자, 필자. 핵심 강점: 로컬 Markdown 저장, 양방향 링크 + 관계 그래프로 노트를 "두 번째 뇌"로 만들 수 있습니다. 데이터는 완전히 당신의 것입니다.
  \item flomo 대상: 짧은 영감은 기록하고 싶은데 긴 글은 쓰기 싫은 사람. 핵심 강점: 카드형 노트. 친구 모멘트 올리듯 생각을 기록하고 \# 태그로 분류하며 타임라인으로 되돌아볼 수 있습니다.
  \item Apple 메모 대상: iPhone/Mac 사용자, 가벼운 기록 수요. 핵심 강점: 빠른 실행, 문서 스캔/드로잉/체크리스트/iCloud 동기화 지원. 무료이며 시스템과 깊게 연동됩니다.
  \item 인샹비지(Evernote) 이런 분께 좋아요: 웹 자료를 대량으로 오려 저장하고 체계적으로 정리하고 싶은 베테랑 사용자. 핵심 장점: 웹 클리퍼가 강력하고, OCR로 이미지 속 글자도 알아봐요. 오래된 만큼 안정적이지만 고급 기능은 유료예요.
\end{itemize}

\booksubsection{PDF 도구}

\begin{itemize}
  \item PDF Expert 대상: iPad/iPhone 사용자, 계약서에 주석이 필요한 비즈니스맨, 학생. 핵심 강점: iOS에서 공인된 최강 PDF 편집기. 매끄러운 주석, 서명, 병합·분할, 양식 작성 지원.
  \item WPS 대상: Word/Excel/PPT/PDF를 동시에 다루는 오피스 사용자. 핵심 강점: 올인원 오피스 스위트. 중국 클라우드/WeChat 파일과 호환성이 좋고, 무료 버전만으로도 일상 충분.
  \item Adobe Acrobat 대상: PDF 규격 준수가 중요한 사람(법률, 디자인, 출판). 핵심 강점: 공식 표준. PDF/A 변환이 정확하고, 크로스 기기 동기화와 전자 서명이 전문적입니다.
\end{itemize}

\booksubsection{스캔(종이 문서 디지털화)}

종이 계약서, 영수증, 숙제, 화이트보드 메모, 신분증 사본… 휴대폰 카메라가 이미 스캐너 대부분을 대체합니다. 핵심 흐름: 선명하게 촬영 자동 크롭·보정 PDF/이미지 내보내기 OCR 선택.

\begin{itemize}
  \item 캠스캐너(CamScanner) 이런 분께 좋아요: 계약서, 영수증, 증명서를 자주 스캔해야 하는 직장인과 학생. 핵심 장점: 가장자리를 자동으로 인식하고 원근을 보정해 주며, OCR로 글자도 추출해요. 여러 페이지를 하나의 PDF로 만들어 클라우드 동기화도 돼요.
  \item Microsoft Lens 대상: 학생, 화이트보드/책 페이지를 Word로 스캔하고 싶은 사람. 핵심 강점: 무료, 광고 없음. OneNote/Word로 원클릭 내보내기, 높은 문자 인식 정확도.
  \item Apple 메모 스캔 대상: iPhone 사용자, 가벼운 스캔 수요. 핵심 강점: 앱 설치 불필요. 메모에서 카메라 아이콘만 누르면 스캔. 자동 그림자 제거, 서명 주석, 무료로 편리합니다.
  \item Baimiao(바이먀오) / CamScanner 웹 버전 / 브랜드 카메라 "문서 스캔": 화웨이, 샤오미, OPPO, vivo 시스템 카메라에는 대부분 "문서/스캔" 모드가 있습니다. 안드로이드는 Google 카메라나 시스템 스캔도 가능. 계약서 한 장 급할 때 유용합니다.
  \item 촬영 팁: 조명이 고르게, 휴대폰은 종이와 평행하게, 네 모서리를 맞추세요. 여러 페이지는 순서대로 찍은 뒤 PDF로 병합합니다. 민감한 신분증은 워터마크와 저장 위치에 신경 쓰고, 사용 후 "최근 삭제된 항목"에서 완전히 정리하세요.
\end{itemize}

\booksection{건강과 운동}

건강 앱은 휴대폰에서 가장 과소평가된 “개인 비서”입니다. 수면, 심박, 운동량을 계속 추적해, 가끔 병원에 가는 것보다 장기 상태를 더 잘 보여 줍니다.

\booksubsection{운동·피트니스}

\begin{itemize}
  \item Keep 대상: 홈트레이닝, 초보자, 체중 감량·체형 교정을 원하는 사람. 핵심 강점: 방대한 무료 강의. HIIT부터 요가까지 폭넓게 제공, AI가 훈련 계획을 자동 생성, 커뮤니티 분위기 활발.
  \item Nike Training Club(NTC) 대상: 전문 코치 지도와 체계적 훈련을 원하는 사람. 핵심 강점: Nike 공식. 전 강의 무료. 동작 시연은 전문 코치가 녹화했고 근육군/기구별로 정리되어 있습니다.
  \item Apple Fitness 대상: Apple Watch 사용자, 활동 링으로 동기부여받고 싶은 사람. 핵심 강점: 매일 세 개의 활동 링(이동/운동/기기기). 목표 시각화, 워치 심박 데이터와 매끄럽게 연결됩니다.
\end{itemize}

\booksubsection{수면 추적}

\begin{itemize}
  \item AutoSleep 대상: Apple Watch 사용자, 매일 밤 깊은/얕은 수면 비율을 알고 싶은 사람. 핵심 강점: 시계를 차고 자면 자동 추적. 수동 시작/정지 불필요. 수면 점수와 취침 조언 제공.
  \item Sleep Cycle 대상: 시계를 차고 자지 않고 휴대폰으로 측정하고 싶은 사람. 핵심 강점: 마이크로 뒤척임/코골이를 분석하고 얕은 수면 단계에서 스마트하게 깨워 줍니다. 기상이 덜 괴롭습니다.
  \item Xiaomi Fitness 대상: 샤오미 밴드/워치 사용자. 핵심 강점: 무료이며 샤오미 밴드와 깊게 연동. 심박/산소포화도/수면 통합 추적, 가성비 우수.
\end{itemize}

\booksubsection{명상·이완}

\begin{itemize}
  \item Calm 대상: 스트레스 많은 직장인, 잠들기 어려운 사람. 핵심 강점: "Daily Calm" 명상 가이드 + 수면 이야기(유명인 낭독). 아름다운 UI, 영어 콘텐츠 풍부.
  \item Tide 대상: 백색 소음으로 집중/수면이 필요한 중국어 사용자. 핵심 강점: 중국산 명품. 자연 사운드 + 명상 코스 + 뽀모도로. 중국어 가이드가 더 친근합니다.
  \item 샤오 슬립(Xiao Sleep) 대상: 수면 유도 사운드와 잠꼬대/코골이 녹음이 필요한 사람. 핵심 강점: 방대한 백색 소음을 중첩 조합 가능(비+선풍기+모닥불). 취침 시 예약 종료. 잠꼬대 녹음은 재미있고 실용적입니다.
\end{itemize}

\booksubsection{여성 건강}

\begin{itemize}
  \item Flo 대상: 생리 기록, 임신 준비, 여성 건강 관리가 필요한 사용자. 핵심 강점: AI가 다음 생리/배란일을 높은 정확도로 예측. 커뮤니티 익명 질문 가능, 섬세한 UI.
  \item 메이유(Meiyou) 대상: 중국 여성 사용자, 중국어 커뮤니티와 건강 정보를 보고 싶은 사람. 핵심 강점: 현지화 콘텐츠 풍부. 생리/임신 준비/임신/육아 단계별 지식 제공.
\end{itemize}

\booksubsection{복약 알림}

\begin{itemize}
  \item 복약 도우미 대상: 장기 복용 만성질환자, 부모님 복약을 챙기는 자녀. 핵심 강점: 바코드로 약 인식, 복약 목록 자동 생성, 시간대별 팝업 알림. 복용 기록을 가족과 공유할 수 있습니다.
\end{itemize}

\booksubsection{데이터 허브}

\begin{itemize}
  \item Apple 건강 / Google Fit 대상: 모든 스마트폰 사용자. 핵심 강점: 휴대폰, 밴드, 타사 건강 앱의 모든 데이터(걸음, 심박, 수면, 생리 등)를 모아 통합 "건강 기록 센터"를 만듭니다.
\end{itemize}

\booksection{여행·이동}

밖으로 나가면 몇 개의 좋은 앱이 무수한 함정을 피해 줍니다. 아래는 여행 전 과정 추천입니다.

\booksubsection{항공·고속철}

\begin{itemize}
  \item UMetrip(항려종횡) 대상: 비행기를 자주 타는 비즈니스맨. 핵심 강점: 공식 민항 데이터. 항공편 동향/게이트 변경/지연 경보가 가장 권위적. 기종과 좌석도 확인 가능.
  \item 12306 대상: 중국 내 기차 여행자 전원. 핵심 강점: 공식 유일 신뢰 예매 채널. 대기 예매(취소표 자동 확보) 지원. 학생표/단체표 발급 편리.
  \item 페이창준(항공편 정보 앱) 이런 분께 좋아요: 항공편 지연에 민감하거나 공항 픽업/배웅을 해야 하는 분. 핵심 장점: 항공편 지연 예측이 정확하고, 일정을 가족에게 공유해 내 위치와 도착 시간을 실시간으로 확인시켜 줄 수 있어요.
\end{itemize}

\booksubsection{지도·내비게이션}

\begin{itemize}
  \item 가오더 지도(Amap) 이런 분께 좋아요: 중국 도시 사용자, 자차/택시 이용자. 핵심 장점: 실시간 교통 상황이 정확하고 혼잡 회피 능력이 뛰어나요. 택시 호출/합승도 한곳에서 해결되고 대중교통 환승 노선도 다양해요.
  \item Baidu Maps 대상: 장소 정보와 실내 내비게이션이 필요한 사용자. 핵심 강점: 파노라마 지도/AR 내비게이션. 관광지/백화점 실내 지도가 세밀하고 음성 팩이 다양합니다.
  \item Google Maps 대상: 해외 여행자. 핵심 강점: 전 세계 커버리지 최강, 현지 리뷰가 현실적. 스트리트 뷰와 오프라인 지도가 타국에서의 구세주입니다.
\end{itemize}

\booksubsection{번역}

\begin{itemize}
  \item Google 번역 대상: 해외 여행자, 다국어 학습자. 핵심 강점: "대화" 실시간 양방향 번역, "카메라" 즉석 단어 번역이 여행 필수품. 오프라인 언어 팩 다운로드 지원.
\end{itemize}

\booksubsection{환율 계산}

\begin{itemize}
  \item 미니멀 환율 이런 분께 좋아요: 해외여행, 해외직구, 국경 간 쇼핑을 하는 분. 핵심 장점: 화면이 정말 깔끔하고 160개 이상의 통화를 지원해요. 환율이 실시간으로 업데이트되고 여러 통화를 한꺼번에 비교할 수 있으며, 숫자를 길게 누르면 빠르게 환산돼요.
\end{itemize}

\booksubsection{날씨}

\begin{itemize}
  \item 기본 날씨 앱 / 중국날씨(China Weather) 이런 분께 좋아요: 모두. 예보와 공식 기상 특보를 확인할 때 첫째로 꼽히는 곳이에요. 핵심 장점: 제조사 날씨 앱과 '중국날씨' 같은 공식 채널은 기상대 특보와 연동되어 있어서, 폭우·태풍·폭염 등 특보가 배너나 알림으로 전송돼요. 권위성이 높아요.
  \item Caiyun Weather(차이윈 날씨) 대상: 분 단위 강수 예보가 필요한 사람. 핵심 강점: "앞으로 두 시간 안에 비가 올까" 정확도가 업계 선두. 강수 구름을 동적 지도로 직관적으로 보여 주고, 비가 오기 전에 먼저 알립니다.
  \item Carrot Weather 대상: 재미와 풍부한 위젯을 원하는 날씨 덕후. 핵심 강점: 독설 AI 음성 날씨 브리핑, 다양한 맞춤 위젯, 전문적인 레이더. 유료 버전이 더 풍부합니다.
  \item 경보 요점: 날씨 앱은 알림을 허용해야 경보를 받을 수 있습니다. 경보는 파랑·노랑·주황·빨강 네 단계이며 빨강이 가장 위험한 재해성 날씨입니다. 안내에 따라 일정을 조정하세요. 더 많은 켜기 방법은 이 장 7절 "날씨 경보 / 자연재해 경보 / 지진 경보"를 참고하세요.
\end{itemize}

\booksubsection{여행 가이드}

\begin{itemize}
  \item 마펑워(Mafengwo) 이런 분께 좋아요: 꼼꼼한 여행 계획을 세우는 걸 좋아하거나 자차 여행·깊이 여행을 즐기는 분. 핵심 장점: 이용자가 직접 쓴 여행기가 풍부하고, 목적지별 즐길 거리·루트·주의점 가이드가 자세해요. 오프라인 가이드를 내려받을 수도 있어요.
  \item 샤오홍슈(RED) 이런 분께 좋아요: 진짜 방문 후기를 보고 숨은 명소·맛집을 찾고 싶은 젊은 여행자. 핵심 장점: 사진과 짧은 동영상으로 직관적이고, 'OO 도시 3일 여행 코스' 같은 콘텐츠가 빠르게 올라와 최신 정보에 강해요. 검색 경험도 좋아요.
\end{itemize}

\booksection{자동화 도구 입문}

자동화는 “한 번 설정으로 오래 편해지는” 지렛대입니다. 매일 반복하는 수십 번의 조작을 0으로 줄여 줍니다. 이 절은 개념 입문만 다루며, 구체적인 활용과 상세 튜토리얼은 23장: 프로그래밍 사고와 자동화 입문을 참고하세요.

\begin{itemize}
  \item iOS "단축어": 애플 공식 자동화 도구로, 프로그래밍 없이 "집에 도착하면 WiFi 켜기 + 밝기 낮추기 + 날씨 읽어주기" 같은 조합 동작을 한 번에 처리할 수 있습니다.
  \item 안드로이드 자동화: 시스템 수준 장면 연동(예: 샤오미 "자동 작업", 화웨이 "스마트 시나리오")이 가장 쉽습니다. Tasker는 기능이 강력하지만 학습 곡선이 가팔라 기크에게 적합합니다.
  \item 핵심 아이디어: 먼저 "매일 반복하는 조작"을 찾고, "하나의 트리거 조건 + 여러 동작"으로 대체할 수 있는지 생각하세요. 상세 예시와 고급 활용은 23장을 참고하세요.
\end{itemize}

\booksection{브라우저와 검색 엔진: 당신의 정보 입구}

브라우저는 인터넷으로 가는 창이고, 검색 엔진은 정보 바다에서 길을 잡아 주는 나침반입니다. 잘 쓰면 우회할 일이 훨씬 줄어듭니다.

\booksubsection{브라우저 작동 원리:}

브라우저에 주소를 입력하면: (1)DNS 서버에 도메인의 IP 주소를 조회하고, (2)그 IP 서버에 HTTP/HTTPS 요청을 보내고, (3)서버가 웹 코드(HTML, CSS, JavaScript)를 반환하며, (4)브라우저가 코드를 해석해 화면으로 렌더링합니다.

\booksubsection{주요 브라우저 추천:}

\begin{itemize}
  \item Microsoft Edge: Chromium 기반으로 빠르고 호환성이 좋습니다. Copilot AI 어시스턴트가 내장되어 웹 요약, 번역, 대화가 가능합니다. 크로스 기기 동기화를 지원해 개인적으로 주력 브라우저로 강력 추천합니다.
  \item Safari: 애플 기기 기본 브라우저. 시스템과 깊게 연동되어 배터리 효율이 높고 프라이버시 보호가 좋지만 확장 기능이 적습니다.
  \item Chrome: 전 세계 사용률 1위, 확장 기능이 풍부합니다. 다만 메모리를 많이 쓰고 모바일에서는 상대적으로 배터리 소모가 큽니다.
\end{itemize}

\booksubsection{검색 엔진 선택:}

\begin{itemize}
  \item Bing: 마이크로소프트 검색 엔진으로 Edge와 깊게 연동됩니다. AI 요약이 링크만이 아니라 답을 직접 제시해 검색 효율이 높습니다. 이미지 검색 체감도 좋습니다.
  \item Yandex: 러시아 최대 검색 엔진. 이미지 검색과 역방향 이미지 검색이 강력합니다. 비슷한 이미지 찾기, 물체 인식에 유용합니다.
  \item Baidu: 중국 내 정보 검색의 주력이지만 광고가 많아 판단력이 필요합니다.
  \item Google: 전 세계에서 가장 정확한 검색 엔진이지만, 중국에서는 특수 네트워크가 있어야 쓸 수 있습니다.
\end{itemize}

\booksubsection{효율적인 정보 검색 가이드:}

\begin{itemize}
  \item AI 검색 활용: Bing의 AI 요약, Edge의 Copilot이 질문에 답을 직접 줘 방대한 결과에서 핵심 정보를 뽑아 줍니다. 링크를 하나씩 여는 것보다 훨씬 효율적입니다.
  \item 허위 정보·광고 식별: 검색 결과의 "광고" 또는 "스폰서" 표시 항목은 유료 프로모션일 뿐 자연 검색 결과가 아닙니다. 출처의 권위를 확인하고 공식 사이트, 유명 매체의 내용을 우선하세요.
  \item 올바른 앱 다운로드: 검색창에 "앱 이름+공식 사이트"를 입력하세요. 예: "WeChat 공식 사이트". 공식 도메인이 표시된 결과를 먼저 클릭하세요. 타사 다운로드 사이트에서 설치하지 마세요. 악성 소프트웨어가 함께 설치되는 경우가 많습니다.
  \item 공식 사이트 식별 팁: Windows 11은 microsoft.com에서, Chrome은 google.com에서 다운로드하세요. 도메인과 브랜드가 일치하는지 확인하는 것이 가장 기본적인 식별 방법입니다.
\end{itemize}

\booksubsection{웹사이트 도메인 지식:}

도메인은 사이트의 “주소”입니다. 읽을 줄 알면 사이트의 성격과 출처를 판단하는 데 도움이 됩니다.

\begin{itemize}
  \item 일반적인 최상위 도메인 의미: .com — 상업 기관, 가장 흔함
  \item .cn — 중국 국가 도메인
  \item .org — 비영리 조직
  \item .net — 네트워크 서비스 기관
  \item .edu — 교육 기관
  \item .gov — 정부 기관
  \item .xyz/.tech/.info — 신흥 범용 도메인. 진입 장벽이 낮아 주의 필요
\end{itemize}

국가/지역 도메인: .jp — 일본

.kr — 한국

.uk — 영국

.de — 독일

.tw — 대만

.hk — 홍콩

보안 팁: .gov와 .edu 사이트는 신청 기준이 높아 보통 더 신뢰할 수 있습니다. .xyz, .top 같은 저가 도메인은 피싱 사이트에 악용되기 쉬우므로 방문 시 주의하세요.

\booksubsection{광고 차단:}

\begin{itemize}
  \item AdGuard 대상: 팝업, 배너 광고, 앱 내 광고로 고통받는 모든 사용자. 핵심 강점: 시스템 수준 광고 차단(브라우저에 국한되지 않음). 앱 내 배너와 스플래시 추적 필터링, 데이터 절약, 페이지 로딩 가속. 유료 버전은 전 기기 지원.
\end{itemize}

\booksubsection{프라이버시 브라우저:}

\begin{itemize}
  \item Brave 대상: 프라이버시를 중시하고 광고 추적을 원치 않는 일상 사용자. 핵심 강점: Chromium 기반. 기본 광고·트래커 차단, Tor 시크릿 모드 내장. Chrome보다 빠르고 메모리 사용량이 낮습니다.
  \item DuckDuckGo 대상: "한 번 검색하면 온통 광고로 쫓기고 싶지 않은" 사용자. 핵심 강점: 동명 검색 엔진 내장(검색 기록 저장 안 함), 원클릭 검색 흔적 삭제("Fire Button"), 크로스 플랫폼 일관 체험.
\end{itemize}

\booksubsection{브라우저 비교 한눈에:}

\begin{center}
\begin{tabularx}{\linewidth}{|X|X|X|}
\hline
\textbf{브라우저} & \textbf{핵심 강점} & \textbf{대상 사용자} \tabularnewline
\hline
Microsoft Edge & Chromium 기반 + Copilot AI 내장. 웹 요약/번역/대화 가능 & 효율을 추구하는 오피스 사용자, AI 집중 사용자(주력 추천) \tabularnewline
Safari & 애플 생태계와 깊게 연동. 배터리 효율 높고 프라이버시 보호 강함 & iPhone/Mac 사용자, 배터리와 시스템 연동 중시 \tabularnewline
Chrome & 전 세계 사용률 1위, 가장 풍부한 확장 생태계 & 확장 기능이 많이 필요한 고급 사용자, 크로스 기기 동기화 사용자 \tabularnewline
Brave & 기본 광고·트래커 차단, 더 빠른 속도 & 프라이버시 중시, 확장 없이도 광고를 막고 싶은 사용자 \tabularnewline
DuckDuckGo & 동명 검색 엔진이 기록을 남기지 않음, 원클릭 흔적 삭제 & 맞춤 추적을 원치 않는 프라이버시 민감 사용자 \tabularnewline
\hline
\end{tabularx}
\end{center}

\begin{tcbpractice}{"취침 전 이완" 자동화 흐름 만들기 연습}
\begin{enumerate}
  \item 트리거 조건 생성: 특정 시간(예: 22:30).
  \item 실행 동작 추가: 음량 낮추기 > 방해 금지 모드 켜기 > 이완 플레이리스트 재생 > 밝기 낮추기.
  \item 저장 후 테스트.
\end{enumerate}

\end{tcbpractice}

\booksection{유용한 시스템 도구: 정리, 스팸 차단, 보안}

생산성·건강 앱 외에 일상 체감을 좌우하는 것은 대개 “방해 덜 받기, 저장 공간 충분히, 사기꾼 못 들어오기” 세 가지입니다. 이 절에서는 자주 쓰는 유틸리티 앱을 정리한 뒤, 휴대폰 기본 스팸 차단 기능을 따로 다룹니다. 휴대폰을 수년 써도 이 스위치들을 한 번도 안 켜 본 사람이 많습니다.

\booksubsection{시스템 정리 및 저장 공간 최적화}

\begin{itemize}
  \item Qingzhuo(징저) 대상: 안드로이드 사용자. 가볍게 캐시, 잔여 파일, 빈 폴더를 정리하고 싶으면서 "휴대폰 관리자" 류의 무거운 앱은 싫은 사람. 핵심 강점: 용량이 작고 광고가 없음. 앱 잔여물 심층 스캔, 대용량 파일 위치 파악, 저장 공간 분석 지원. 상업용 "원클릭 가속"보다 절제되어 있어 손재주 있는 사용자의 정기 유지보수에 적합. 주의: 정리 전에 경로를 확인하세요. WeChat, 갤러리, 문서 폴더는 무턱대고 전체 선택하지 마세요. 시스템 캐시는 제조사 기본 정리 도구를 우선하고, 타사 도구는 보조로만 쓰세요.
  \item 시스템 기본 저장 공간 정리 대상: 모든 사용자, 특히 부모님과 설정이 부담스러운 분. 핵심 강점: 안드로이드 [설정] > [저장 공간] > [정리 및 가속], iOS [설정] > [일반] > [iPhone 저장 공간]. 시스템이 자사 앱 캐시 구조를 더 잘 알아 오삭제 위험이 가장 낮습니다. 1순위로 선택하세요.
  \item WeChat/QQ 내장 저장 관리 대상: 채팅이 많은 사용자. 핵심 강점: WeChat [나] > [설정] > [일반] > [저장 공간]에서 채팅, 파일 유형별로 세밀하게 정리할 수 있습니다. 외부 "WeChat 정리"류 도구보다 안전합니다.
  \item 원칙: 시스템/앱에 있는 기능이면 타사를 쓰지 말 것. 경로/유형별로 고를 수 있으면 "원클릭 전체 삭제"를 쓰지 말 것. 정리의 목표는 공간 확보이지 중요한 데이터를 함께 날리는 것이 아닙니다.
\end{itemize}

\booksubsection{스팸 차단·보안 유틸리티}

\begin{itemize}
  \item 국가 사기 방지 센터(National Anti-Fraud Center) 이런 분께 좋아요: 모두. 특히 어르신, 재무 관련 직장인, 온라인 쇼핑을 자주 하거나 낯선 전화를 자주 받는 분. 핵심 장점: 공식 앱으로, 수신 전화 경고, 문자 위험 알림, 의심스러운 웹사이트 검사 기능을 켤 수 있어요. 공안 사기 방지 데이터베이스와 연동되어 '3중 방어선' 중 공적 신뢰가 가장 강한 한 축이에요.
  \item 브랜드 보안 센터 / 휴대폰 관리자 대상: 해당 브랜드 사용자. 핵심 강점: 화웨이 휴대폰 관리자, 샤오미 보안 센터, OPPO 휴대폰 관리자, vivo 휴대폰 관리자 등. 스팸 차단, 바이러스 검사, 권한 관리, 결제 보호를 통합합니다. 보통 무료이며 시스템 하위 레벨과 협업이 더 좋습니다. 두 번째 상업용 보안 소프트웨어를 깔아 권한 충돌을 만들지 마세요.
  \item 통신사 스팸 차단 서비스 대상: 고빈도 발신, AI 음성 전화 폭탄에 시달리는 사용자. 핵심 강점: China Mobile은 KTFSR을 10086으로 전송해 개통, China Telecom은 "Tianyi Anti-Spam" 팔로우, China Unicom은 "Unicom Mobile Manager"에서 설정. 차단이 통신사 측에서 이루어져 휴대폰이 울리지 않아도 스팸 번호를 막아 줍니다.
  \item AdGuard 대상: 팝업, 배너 광고, 앱 내 광고로 고통받는 모든 사용자. 핵심 강점: 시스템 수준 광고 차단(브라우저에 국한되지 않음). 앱 내 배너와 스플래시 추적 필터링, 데이터 절약, 페이지 로딩 가속. 유료 버전은 전 기기 지원.
  \item "정리 가속 / 스팸 차단" 류 올인원 패키지 주의: 일부 타사 관리자 앱은 백그라운드 상시 실행, 연락처·문자 권한 반복 요청. 차단 능력이 시스템보다 낫지도 않은데 배터리와 프라이버시 위험만 키웁니다. 시스템 기본 기능을 먼저 켜고, 그다음 공식 사기 방지 앱을 고려하세요.
\end{itemize}

\textbf{유틸리티 앱 선택 원칙}\par

\begin{itemize}
  \item 시스템 우선: 차단, 정리, 권한 관리는 휴대폰 기본 스위치를 먼저 켜세요.
  \item 공식 보완: 국가 사기 신고 센터, 통신사 스팸 차단은 공식 채널이므로 함께 쓸 수 있습니다.
  \item 타사 절제: 같은 기능은 하나만 설치하세요. 권한 충돌과 백그라운드 자원 경쟁을 피하세요.
  \item 적을수록 많다: 도구가 많을수록 알림과 권한 노출면이 커지고, "도구 자체"에게 방해받을 확률이 올라갑니다.
\end{itemize}

\booksubsection{휴대폰 기본 스팸 차단 기능: 브랜드별 켜기}

아래는 브랜드별 진입 경로입니다. 이름은 시스템 버전에 따라 조금 다를 수 있으니 설정에서 “스팸”, “차단”, “알 수 없는 발신”으로 검색해 보세요. 8장 “스팸 전화 대응”과 서로 보완합니다. 여기는 스위치 위치에 초점을 맞추고, 그쪽은 번호 표시와 오표시 이의신청을 다룹니다.

\begin{center}
\begin{tabularx}{\linewidth}{|X|X|X|}
\hline
\textbf{플랫폼} & \textbf{진입 경로} & \textbf{권장 켜기 항목} \tabularnewline
\hline
iPhone / iOS & {}[설정] > [전화] / [메시지] & 알 수 없는 발신자 음소거, 차단된 연락처, 메시지의 "알 수 없는 발신자 필터" \tabularnewline
Xiaomi / Redmi & 전화 앱 > 설정 > 스팸 차단 또는 보안 센터 & 스마트 차단, 모르는 번호 한 번 울림 끊김, 문자 키워드 필터 \tabularnewline
Huawei / Honor & 전화 > 스팸 차단, 휴대폰 관리자 > 스팸 차단 & 발신자 식별, 블랙리스트, 문자 스마트 차단 \tabularnewline
OPPO / realme / OnePlus & 전화 앱 > 스팸 차단 또는 기기 > 스팸 차단 & 수신 차단, 보이스피싱 식별, 문자 차단 규칙 \tabularnewline
vivo / iQOO & 전화 > 스팸 전화 차단, i Manager & 스팸 차단, 보이스피싱 경보, 메시지 필터 \tabularnewline
Samsung & 전화 > 설정 > 번호 차단 / 발신자 표시 및 스팸 전화 & 알 수 없는 번호 차단, 스팸 전화 식별(지역 지원도 상이) \tabularnewline
\hline
\end{tabularx}
\end{center}

\begin{tcbinfo}{3중 방어 조합 권장}
일상에는 “시스템 차단 + 통신사 스팸 차단 + 국가 사기 신고 센터” 3중 조합을 권장합니다. 시스템은 표시와 블랙리스트를, 통신사는 네트워크 측에서 고빈도 발신을, 사기 신고 센터는 공식 경보를 맡습니다. 셋 다 최고 강도로 켤 필요는 없습니다. 택배, 배달, 병원 전화까지 오차단하지 않도록 연락처와 최근 연락처는 통과시키세요.

\end{tcbinfo}

\begin{tcbpractice}{실습: 오늘 스팸 차단 켜기}
\begin{enumerate}
  \item 전화 앱 설정에서 "스팸 차단 / 알 수 없는 발신자 음소거"를 찾아 스마트 차단을 켭니다.
  \item 부모님 휴대폰도 한 번 확인하세요. 같은 경로 + 통신사 스팸 차단 개통(China Mobile은 KTFSR을 10086으로 전송).
  \item 앱 스토어에서 "국가 사기 신고 센터"를 설치하고, 가입 후 수신 경보를 켭니다.
  \item 이미 표시된 테스트 번호로 전화를 걸거나(또는 차단 기록을 보고) 기능이 작동하는지 확인합니다.
\end{enumerate}

\end{tcbpractice}

\booksection{날씨 경보, 자연재해 경보, 지진 경보}

스팸 차단이 막는 것은 “전화 사기”이고, 경보 기능이 막는 것은 “자연재해와 극단적 날씨”입니다. 날씨 앱만 깔고 온도만 보는 사람이 많은데, 알림과 지진 경보는 안 켠 경우가 많습니다. 폭우 적색 경보, 태풍, 지진이 오기 직전 수 초\textasciitilde{}수십 초 동안 휴대폰은 소중한 대피 시간을 벌어 줄 수 있습니다.

\booksubsection{날씨 경보: 폭우, 태풍, 폭염 받는 법}

\begin{itemize}
  \item 공식 채널 우선: 중국기상국 / 중앙기상대 / 각 성 기상국이 발표하는 특보가 가장 권위 있어요. '중국날씨' 위챗 공식 계정이나 웨이보를 팔로우하거나, '중국날씨' 같은 공식 날씨 앱을 내려받으세요.
  \item 시스템 날씨 앱: 화웨이 날씨, 샤오미 날씨, OPPO 날씨, vivo 날씨, 애플 "날씨" 등 대부분 현지 기상대 경보를 푸시합니다. 날씨 앱 설정에서 "기상 경보 / 재해성 날씨" 알림이 켜져 있는지 확인하세요.
  \item Caiyun Weather 등 분 단위 강수: 출퇴근, 아이 등원 전에 "이 30분 우산을 챙길까" 판단하기에 좋습니다. 경보 능력은 여전히 공식 기상대가 기준입니다.
  \item 경보 색상 의미(일반): 파랑: 주의. 평소처럼 이동하되 업데이트를 확인.
  \item 노랑: 경계 강화. 불필요한 야외 활동을 줄이세요.
  \item 주황: 심각. 실내에 머물고 위험 지역에서 멀리하세요.
  \item 빨강: 매우 심각. 휴교·휴업·대피 안내를 따르세요. 침수 도로, 산사태, 강풍 구간을 무리하게 지나가지 마세요.
\end{itemize}

일반 경보 유형: 폭우, 폭설, 태풍, 폭염, 한파, 강풍, 뇌우, 우박, 도로 결빙, 짙은 안개, 가뭄, 지질재해 위험.

\booksubsection{자연재해 경보: “비”만이 아닙니다}

\begin{itemize}
  \item 발표 주체: 기상대는 기상 관련을 담당하고, 지질재해·산사태 등은 자연자원·수리·응급 관리 부처가 공동 발표합니다. 지방 정부도 비상 방송, 문자, 12379 등 채널로 발표합니다.
  \item 다중 채널 수신 방법: 휴대폰 시스템 알림 + 날씨 앱 경보.
  \item WeChat에서 현지 "응급관리국""기상국" 공식 계정 팔로우.
  \item 12379 국가 경보 정보 발신 채널(문자/서비스, 지역 개통 여부에 따름).
  \item TV, 비상 방송, 마을/커뮤니티 확성기 — 산간, 홍수기에는 휴대폰만 보지 마세요.
\end{itemize}

야외·이동 시 주의: 태풍 시 광고판 근처와 나무 아래를 피하고, 폭우 시 지하 주차장·터널·하부 통과 교량에 들어가지 마세요. 산간 여행 중 "지질재해 기상 위험" 경보가 있으면 즉시 비탈면·계곡에서 벗어나세요.

응급 안내서 연동: 재해 후 연락 두절, 배터리 방전, 위치 구조 요청이 필요하면 32장 지진/재해 장면과 7장 위성 통신·무통화 섹션을 참고하세요.

\booksubsection{지진 경보: “지진 예보”와 다릅니다}

지진 경보(EEW)는 지진파(특히 파괴력이 큰 S파)가 목표 지역에 도착하기 전에, 전자기파가 지진파보다 빠르고 관측소가 P파를 더 빨리 포착하는 원리를 이용해 수 초\textasciitilde{}수십 초 전에 경보를 보냅니다. 지진 발생 여부를 예측하지는 못하지만, “가까운 곳 대피, 화원 차단, 엘리베이터 탈출”의 시간 벌어 줍니다. 진앙에 가까울수록 경보 시간이 짧아 심지어 채 오지 않을 수 있습니다.

\begin{itemize}
  \item 공식·업계 방안: 중국 대륙에서는 중국 지진 조기경보망 관련 기능과 청두 고신 재해감소연구소 등 기관이 추진하는 경보 서비스가 널리 쓰여요. 일부 제조사 휴대폰은 미리 설치되어 있거나 네트워크에 연결해 경보를 받을 수 있어요.
  \item 타사 앱: 앱 스토어에서 "지진 경보" 등 공식/기관 관련 앱을 검색하세요. 설치 전에 개발자를 확인하고, 무관한 권한을 대거 요구하는 짝퉁 "가짜 경보" 앱을 피하세요.
  \item 경보를 받은 후: 실내라면 가까운 곳에 엎드리기(책상 아래, 내력벽 모서리, 견고한 가구 옆)하고 머리·목을 보호하세요. 창문, 외벽, 높은 가구에서 멀리하세요. 무작정 계단으로 뛰거나 뛰어내리지 마세요. 엘리베이터에 있다면 가장 가까운 층에서 즉시 내리세요. 야외라면 전선, 광고판, 위험 건물에서 떨어진 넓은 곳으로 가세요. 진동 후 가스와 전기를 확인하고 여진에 주의하세요.
  \item 알림을 끄지 마세요: 지진 경보는 "무음/방해 금지 돌파"가 허용되어야 합니다. 잠들기 전에 손이 닿는 곳에 휴대폰을 두고 벨소리를 0으로 두지 마세요.
\end{itemize}

\booksubsection{휴대폰 기본 지진 경보: 브랜드별 켜기}

아래 경로는 시스템 버전에 따라 다를 수 있으니 설정에서 “지진”, “경보”, “안전 경보”로 검색하세요. iOS 자체의 중국 지진 경보 지원은 제한적이라 보통 신뢰할 수 있는 타사 경보 앱을 설치하고 위치·알림 권한을 유지해야 합니다.

\begin{center}
\begin{tabularx}{\linewidth}{|X|X|X|}
\hline
\textbf{플랫폼} & \textbf{일반 경로} & \textbf{설명} \tabularnewline
\hline
Xiaomi / Redmi & 휴대폰 관리자 / 보안 센터 > 지진 경보 또는 설정에서 "지진" 검색 & 일부 버전은 시스템에 경보 서비스 내장. 네트워크 연결과 위치 권한 필요 \tabularnewline
Huawei / Honor & 설정 > 보안 > 비상 경보 알림 / 지진 경보 & 이름이 "비상 경보 알림"일 수 있으며 지진 등이 포함됩니다 \tabularnewline
OPPO / realme / OnePlus & 설정 > 보안(SOS 긴급 연락) > 지진 경보 & 일부 기종은 지역 서비스 커버리지와 관련 \tabularnewline
vivo / iQOO & 날씨 앱 또는 설정 > 보안/기타 > 지진 경보 & 실제 시스템 버전 메뉴 기준 \tabularnewline
iPhone / iOS & 시스템 설정에 완전한 내장 스위치가 없을 수 있음. 신뢰할 수 있는 지진 경보 앱을 설치하고 알림을 허용하세요 & 앱에 백그라운드 실행·위치 권한 필요. 시스템 방해 금지가 경보를 차단하지 않도록 하세요 \tabularnewline
Samsung 등 & 설정 > 안전 및 긴급 상황 / 무선 긴급 경보 & "긴급 경보" 범위는 지역·통신사에 따라 다릅니다 \tabularnewline
\hline
\end{tabularx}
\end{center}

\begin{tcbwarn}{지진 경보에 대한 흔한 오해}
\begin{itemize}
  \item 경보 $\neq$ 예보: "내일 지진이 날까"를 알려 주지 못합니다. 이미 발생한 지진의 파동이 아직 도착하지 않았을 때만 경보합니다.
  \item 어디서든 수십 초를 벌 수 있는 것은 아닙니다: 진앙 근처는 수 초뿐이거나 아예 채 오지 않을 수 있으며, 경보는 주로 주변 지역에 이로움을 줍니다.
  \item 알림을 끄면 없는 것과 같습니다: 아무리 좋은 경보라도 무음, 백그라운드 정리, 위치 거부를 당하면 받지 못할 수 있습니다.
  \item 짝퉁 앱: 연락처·문자 권한을 요구하는 "지진 예보" 소프트웨어를 경계하세요. 시스템 내장 또는 정규 기관 솔루션을 우선하세요.
\end{itemize}

\end{tcbwarn}

\begin{tcbinfo}{경보 능력 조합}
일상 권장: “시스템 날씨 경보 + 공식 기상/응급 채널 + 지진 경보(시스템 또는 정규 앱)”. 스팸 차단과 마찬가지로 같은 류의 소프트웨어를 잔뜩 깔 필요는 없습니다. 핵심은 알림 권한 허용, 위치 허용, 방해 금지 돌파입니다. 부모님 휴대폰의 스팸 차단을 켰다면 날씨 경보와 지진 경보도 함께 확인하세요.

\end{tcbinfo}

\begin{tcbpractice}{실습: 10분 만에 경보 설정 점검}
\begin{enumerate}
  \item 시스템 날씨 앱을 열고 설정에서 "기상 경보 / 재해성 날씨" 알림을 켭니다.
  \item 시스템 설정에서 "지진" 또는 "경보"를 검색해 지진 경보/비상 경보를 켭니다(안드로이드). iPhone은 정규 지진 경보 앱을 설치하고 알림을 허용하세요.
  \item 현지 기상국 또는 응급관리국의 공식 계정/서비스 계정을 팔로우하고 푸시를 허용하세요.
  \item 경보 소리가 방해 금지로 완전히 무음 처리되지 않는지 확인하세요. 부모님 휴대폰도 같은 절차로 확인하세요.
\end{enumerate}

\end{tcbpractice}

\booksection{크로스 기기 협업: 원격 제어, 고속 전송, 계정 보안}

휴대폰이 편한 이유는 대개 PC·태블릿과 “팀”을 이루기 때문입니다. 이 절에서는 자주 쓰는 도구를 보충합니다. 휴대폰 원격 PC 제어, AI 오피스·원격 협업, 동적 배경 화면, LAN 고속 전송, 크로스 기기 클립보드, OTP 2FA, 종이 스캔과 일정 연결 등. 보다 체계적인 다중 기기 연동은 17장을 참고하세요.

\booksubsection{휴대폰으로 PC 원격 제어: Sunlogin / ToDesk}

\begin{itemize}
  \item Sunlogin 원격 제어 대상: 외출 중 집/회사의 Windows나 Mac을 조작해야 하는 사용자. 부모님 휴대폰을 원격으로 수리하는 데도 자주 쓰입니다(27장 참고). 핵심 강점: 중국산 생태계가 성숙. 휴대폰 앱으로 PC 데스크톱 조작, 파일 전송, 전원 켜기/끄기(일부는 하드웨어 필요). 무료 버전이 개인 가벼운 사용에 충분합니다.
  \item ToDesk 대상: 연결 안정성과 저지연 원격 업무/운영을 추구하는 사용자. 핵심 강점: 기기 코드 + 비밀번호 연결. 휴대폰에서 PC 제어, PC에서 휴대폰 제어(버전에 따름), 파일 전송, 다중 모니터 전환이 매끄럽습니다. 중국 내 노드 체감이 보통 좋습니다.
  \item 기본 사용법: PC에 클라이언트를 설치하고 로그인/온라인 유지 휴대폰에 앱 설치 기기 식별 코드 + 연결 비밀번호(또는 동일 계정)로 세션 생성 휴대폰 트랙패드/제스처로 데스크톱 조작.
  \item 보안 금지선(꼭 기억): 모르는 사람이 "고객센터/공안·검찰·법원/통신사"를 사칭하며 Sunlogin, ToDesk를 설치하게 하고 휴대폰이나 PC를 원격 제어하게 하려 하면 모두 거절하세요. 전형적인 사기 수법입니다(27장 참고).
  \item 직접 쓸 때: 2차 인증 켜기, 기본 비밀번호 변경, 쓰지 않을 때 로그아웃 또는 무인 모드 끄기. 회사 PC는 사내 보안 규정을 따르세요.
  \item 공공장소에서 원격 업무를 할 때는 주변에 화면이나 비밀번호를 훔쳐보는 사람이 없는지 확인하세요.
\end{itemize}

\booksubsection{AI 오피스 소프트웨어: Trae / WorkBuddy / Dumate}

\begin{itemize}
  \item Trae 대상: 프로그래밍·문서 협업 사용자, 중국어 환경에 친화적. 핵심 강점: ByteDance가 내놓은 AI IDE/코딩 어시스턴트. SOLO 모드는 한 문장으로 프로젝트를 생성합니다. 일상 코딩·문서 수정 효율이 높습니다. PC에서 실행하면 모바일에서 원격 제어 또는 공식 연동 기능으로 작업을 확인·추적할 수 있습니다(제품 최신 형태 기준). 자세한 내용은 41장 AI 코딩 도구 전경을 참고하세요.
  \item WorkBuddy 대상: 회의록, 이메일, 보고서, 사무 문서를 AI가 정리해 주길 원하는 직장인. 핵심 강점: 오피스 장면을 겨냥한 AI 어시스턴트. "초안—윤색—요약" 같은 고빈도 문서 작업에 적합합니다. 일부 워크플로는 휴대폰에서 시작하고 무거운 일은 PC에서 처리할 수 있습니다.
  \item Dumate 대상: 다인 협업, 문서·지식 관리 통합이 필요한 팀/개인. 핵심 강점: AI 협업형 오피스 도구. 자료 모으기와 공동 편집을 강조합니다. 휴대폰에서는 조회와 가벼운 편집이 가능하고, 복잡한 조판은 PC 대화면에서 처리하는 것이 좋습니다.
  \item "휴대폰으로 PC 원격 제어"의 의미: 많은 AI 오피스 기능은 PC/웹에서 돌아갑니다. 휴대폰에서 Sunlogin/ToDesk로 PC에 연결하거나 제품 자체의 모바일 앱/웹을 쓰면 출퇴근 중에도 PC의 AI 작업을 "리모컨"할 수 있습니다. 원칙: 무거운 계산과 긴 문서는 PC에, 진행 상황 확인과 몇 줄 수정은 휴대폰에.
  \item 프라이버시 주의: 계약서, 신분증, 고객 데이터를 업로드하기 전에 제품 개인정보 처리 방침과 기업 컴플라이언스 요구를 확인하세요. 회사 자료는 사내 지정 도구를 우선하세요.
\end{itemize}

\booksubsection{인공 데스크톱: miHoYo 동적 배경 화면(휴대폰 + PC)}

\begin{itemize}
  \item Artificial Desktop(인공 데스크톱) 대상: miHoYo 캐릭터를 좋아하고 무료 동적 배경 화면을 원하는 사용자. 핵심 강점: miHoYo가 내놓은 무료 동적 배경 화면 앱. 휴대폰과 PC 지원. 캐릭터 Live2D/동적 연출을 데스크톱으로 쓸 수 있으며, 일부 기종/시스템에서는 잠금 화면 또는 배경 화면으로 설정할 수 있습니다.
  \item 배터리·성능: 동적 배경은 정지 이미지보다 배터리를 더 쓰고 GPU를 점유합니다. 배터리가 낮거나 오래된 기종은 충전 중에만 쓰거나 프레임 속도를 낮추세요(앱이 옵션을 제공한다면).
  \item 설치 출처: 공식 사이트/정규 앱 스토어를 우선하고, 타사 "크랙" 버전에 악성 코드가 섞이지 않도록 주의하세요. PC 버전은 설치 경로와 시작 프로그램 등록 여부를 확인하세요.
\end{itemize}

\booksubsection{같은 네트워크에서 파일 빠르게 보내기: LocalSend / 페이거추안수(Feige) / 에어드롭(AirDrop)}

\begin{itemize}
  \item LocalSend 대상: 안드로이드/Windows/iOS/Mac/Linux를 혼용하고, 퍼블릭 클라우드 없이 크로스 플랫폼 고속 전송을 원하는 사용자. 핵심 강점: 오픈소스, LAN 직접 연결. 같은 Wi-Fi에서 사진·영상·문서를 상호 전송. 데이터 케이블 불필요, 제조사 서버 경유 없음. 속도는 공유기에 달렸습니다.
  \item 페이거추안수(Feige) 이런 분께 좋아요: 사무실 Windows PC끼리, 예전 방식의 사내 네트워크 환경. 핵심 장점: 오래된 LAN 전송/채팅 도구로 사내망에서 큰 파일을 주고받기에 좋아요. 반드시 정식 경로에서 내려받고, 회사 내부망에서는 IT 정책을 지켜야 해요.
  \item AirDrop 대상: iPhone, iPad, Mac 사용자. 핵심 강점: 애플 생태계 원클릭 공유. 블루투스 검색 + Wi-Fi 직접 전송으로 데이터가 절약됩니다. "연락처만"으로 설정해 모르는 사람의 방해를 줄일 수 있습니다.
  \item 기타 보충: Xiaomi Share, Huawei Share, vivo Share 등 브랜드 솔루션은 같은 브랜드 기기 사이에서 빠릅니다. WeChat "파일 전송 도우미"는 편하지만 파일 크기와 로그인 기기 제한이 있어, 대용량 파일에는 LocalSend/AirDrop을 더 권장합니다.
  \item 사용 팁: 양쪽 모두 같은 Wi-Fi에 연결하세요. 공용 Wi-Fi에서 민감 파일 전송은 위험하니 개인 핫스팟을 우선하세요. 전송 후 수신 측에서 파일이 완전한지 확인하세요.
\end{itemize}

\booksubsection{크로스 기기 클립보드 동기화: KDE Connect / 휴대폰-PC 상호 전송}

\begin{itemize}
  \item Apple 유니버설 클립보드 대상: iPhone + Mac 사용자. 핵심 강점: 같은 Apple ID로 로그인하고 Handoff를 켜면 휴대폰에서 복사 Mac에서 바로 붙여넣기(그 반대도 가능). 텍스트와 일부 파일 모두 가능합니다.
  \item KDE Connect 대상: 안드로이드 + Windows/Linux/Mac. 오픈소스 클립보드 동기화, 알림 미러링, 원격 입력을 원하는 사용자. 핵심 강점: LAN 페어링 후 클립보드 동기화, 휴대폰 알림을 PC로 푸시, 휴대폰을 슬라이드 넘기기/트랙패드로 사용. 데이터가 타사 계정을 거치지 않습니다.
  \item 브랜드 "휴대폰-PC 상호 전송": 화웨이 다중 화면 협업/슈퍼 터미널, 샤오미 미아오샹, vivo Share, OPPO Connect, Honor Connect 등. 크로스 화면 클립보드와 파일 드래그를 자주 지원하며, 같은 브랜드에서 체감이 가장 좋습니다.
  \item Windows + 안드로이드 공식 채널: "Phone Link / Link to Windows"(Your Phone의 후속)로 PC에서 휴대폰 알림과 사진을 보고, 지원되는 장면에서 복사·붙여넣기를 할 수 있습니다. Samsung 등 일부 기종은 더 깊게 통합됩니다.
  \item 보안 주의: 클립보드에 비밀번호, 인증번호가 남을 수 있습니다. 공용 PC에서는 전역 동기화를 신중히 쓰고, 사용 후 비밀번호 입력창에 다시 붙여넣어 유출되는 일을 막으세요. 일부 입력기도 클립보드를 읽으니 권한에 주의하세요.
\end{itemize}

\booksubsection{OTP 2FA: Aegis / Microsoft Authenticator}

\begin{itemize}
  \item 인증기 앱을 문자 대신 쓰는 이유: 문자 인증번호는 SIM 재발급 납치에 노출될 수 있습니다. TOTP OTP는 로컬에서 시간 알고리즘으로 생성되어 휴대폰 신호에 의존하지 않으며, 피싱·차단에 더 강합니다(악성 코드와 스크린샷 촬영은 여전히 조심해야 합니다).
  \item Aegis Authenticator 대상: 안드로이드 사용자, 오픈소스·로컬 제어 선호. 핵심 강점: 오픈소스 무료. 암호화 백업과 가져오기·내보내기 지원. 깔끔한 UI, 지문/얼굴로 토큰 목록 잠금 해제.
  \item Microsoft Authenticator 대상: Microsoft 계정, Office 365/기업 계정 사용자. iOS/안드로이드 크로스 플랫폼. 핵심 강점: 푸시 승인 로그인, 비밀번호 모니터링, 계정 백업(제품 버전에 따름). Windows/기업 신원 체계와 긴밀히 결합됩니다.
  \item Google Authenticator / 2FAS 등: 생태계 안에서도 흔히 쓰입니다. 핵심은 암호화 백업이 있어 기기를 바꿔도 토큰을 잃지 않는 것입니다.
  \item 기기 변경 시 필수: 기존 기기에서 이전하거나 서비스의 "백업 코드/복구 키"를 사용하세요. 백업 코드는 종이에 인쇄해 두 곳에 보관하고, 같은 휴대폰에만 저장하지 마세요.
  \item 비밀번호 관리자와: 일부 비밀번호 관리자도 TOTP를 저장할 수 있어 편하지만 위험이 모입니다. 고가치 계정은 "비밀번호 관리자 + 별도 인증기"로 역할을 나누는 것이 좋습니다. 자세한 내용은 14장을 참고하세요.
\end{itemize}

\booksubsection{추천 조합: 하루 동안 엮는 법}

\textbf{실용 워크플로 예시}\par

\begin{itemize}
  \item 아침: TickTick에서 오늘 일정 확인 날씨/경보 훑기 PC에서 Trae/문서 열고 휴대폰은 옆에 두기.
  \item 낮: 종이 자료는 스캔 앱으로 PDF 변환 후 LocalSend/AirDrop으로 PC에 전송. 긴 문서는 PC에서 AI 오피스 도구로 초안 작성.
  \item 외출: ToDesk/Sunlogin으로 PC에 연결해 작업 계속. 중요한 계정 로그인은 Aegis/Microsoft Authenticator OTP 사용.
  \item 저녁: TickTick에서 완료 항목 체크. 필요하면 동적 배경 화면으로 분위기 전환.
\end{itemize}

\begin{tcbwarn}{원격 제어·계정 보안 수칙}
\begin{itemize}
  \item 남이 원격 제어 소프트웨어를 설치하라고 하면 전화를 끊고 직접 공식 고객센터에 연락하세요.
  \item 인증번호, OTP는 "택배""고객센터"를 포함해 누구에게도 알려 주지 마세요.
  \item 회사 PC 원격 제어나 AI 소프트웨어 설치 전에 먼저 IT에 사용 가능한지 물어보세요.
  \item 공용 네트워크에서 파일 전송이나 원격 로그인을 할 때는 특히 조심하세요.
\end{itemize}

\end{tcbwarn}

\begin{tcbpractice}{실습: "크로스 기기" 연결 한 번 뚫어 보기}
\begin{enumerate}
  \item 휴대폰과 PC를 같은 Wi-Fi에 연결하고 LocalSend(또는 AirDrop)로 사진을 주고받아 보세요.
  \item TickTick에 "오늘 밤 8시: 스캔 파일 정리" 작업을 만들고 휴대폰과 PC에서 모두 보이는지 확인하세요.
  \item 이메일이나 WeChat에 2FA를 켜고(지원한다면) Aegis 또는 Microsoft Authenticator로 연결한 뒤 백업 코드를 저장하세요.
  \item 자신의 PC에 Sunlogin 또는 ToDesk를 설치하고 자신의 휴대폰으로 한 번 성공적으로 연결한 뒤 로그아웃하거나 비밀번호를 강화하세요. 흐름은 익히되 보안 구멍은 남기지 마세요.
\end{enumerate}

\end{tcbpractice}
